\section{Lie groups: exponential coordinates, the local--global distinction, and nilpotent group laws}\label{sec:liegroups}
\subsection{Matrix groups}
For a matrix Lie group $G\subseteq\GL_d$, its Lie algebra $\g$ is the
tangent space at the identity. The left-invariant vector field associated
with $X\in\g$ is $X^L(g)=gX$, and the right-invariant field is $X^R(g)=Xg$. Regarding these as vector fields on the open
subset $\GL_d$ of the vector space of matrices, the bracket of two vector
fields $U,V$ is $[U,V](g)=DV(g)[U(g)]-DU(g)[V(g)]$, and since
$DX^L(g)[H]=HX$,
\[
 [X^L,Y^L](g)=(gX)Y-(gY)X=g(XY-YX)=[X,Y]^L(g).
\]
Thus the Lie bracket of $\g$, defined through left-invariant fields, is the
matrix commutator. The integral curve of $X^L$ through $I$ solves
$g'=gX$, $g(0)=I$; the matrix exponential $e^{tX}$ solves this equation, and
the solution is unique (differentiate $g(t)e^{-tX}$), so the Lie-group
exponential is the ordinary matrix exponential. This proves the matrix
statements of the source page.

For matrix Lie algebras, and more generally for closed Lie subalgebras of a
Banach algebra, the fact that the BCH logarithm lies in the Lie algebra can
be obtained directly from the differential equation of \cref{thm:ode},
without the formal theory of \cref{sec:hopf}.
\begin{proposition}[The BCH logarithm in a closed Lie subalgebra]\label{prop:closedlie}
Let $\cA$ be a real or complex unital Banach algebra. There is $\delta>0$
such that for every closed linear subspace $\g\subseteq\cA$ that is closed
under the commutator $[a,b]=ab-ba$, and all $X,Y\in\g$ with
$\norm X+\norm Y<\delta$, the logarithm $Z=\log(e^Xe^Y)$ of
\cref{cor:bchlog}, that is, the sum of the BCH series, lies in $\g$. In
particular, for a matrix Lie algebra $\g\subseteq\mathfrak{gl}_d$ and $X,Y\in\g$
sufficiently small, $e^Xe^Y=e^Z$ with $Z\in\g$.
\end{proposition}
\begin{proof}
Take $\delta$ as in \cref{thm:ode}, shrunk so that $Z(t)=\log(e^Xe^{tY})$
satisfies $\norm{Z(t)}\leq\tfrac18$ for $|t|\leq1$; then $Z(0)=X$ and
$Z'(t)=F(Z(t))$ with $F(W)=\beta(\ad_W)Y$. The map $F$ sends $\g$ into
itself: $\ad_W^nY$ is an iterated commutator of elements of $\g$, and $\g$
is closed, so the series $\sum_n\bplus n\ad_W^nY$ lies in $\g$. Moreover $F$
is Lipschitz on the ball $\norm W\leq\tfrac38$: with
$\norm{\ad_W}\leq2\norm W\leq\tfrac34$, one has
$\norm{\ad_W^{n}-\ad_{W'}^{n}}\leq n(\tfrac34)^{n-1}\norm{\ad_W-\ad_{W'}}
\leq2n(\tfrac34)^{n-1}\norm{W-W'}$, and $|\bplus n|\leq1$ by
\eqref{eq:bernrec}, so $\norm{F(W)-F(W')}\leq K\norm Y\,\norm{W-W'}$ with
$K=2\sum_{n\geq1}n(\tfrac34)^{n-1}$.

Let $d(t)=\operatorname{dist}(Z(t),\g)$, a continuous function on $[0,1]$
with $d(0)=0$. For $w\in\g$ with $\norm{Z(t)-w}$ close to $d(t)$ (so that
$\norm w\leq\tfrac38$) and $h>0$, the element $w+hF(w)$ lies in $\g$, hence
\begin{align*}
 d(t+h)&\leq\norm{Z(t+h)-w-hF(w)}\\
 &\leq\norm{Z(t+h)-Z(t)-hF(Z(t))}+\norm{Z(t)-w}+h\norm{F(Z(t))-F(w)}\\
 &\leq o(h)+(1+hK\norm Y)\norm{Z(t)-w}.
\end{align*}
Taking the infimum over $w$ gives $d(t+h)\leq(1+hK\norm Y)d(t)+o(h)$, so
the upper right Dini derivative of $d$ is at most $K\norm Y\,d(t)$.
Gronwall's inequality yields $d\equiv0$ on $[0,1]$, and since $\g$ is
closed, $Z=Z(1)\in\g$. A matrix Lie algebra is finite-dimensional, hence
closed, and closed under the commutator, so the last assertion follows.
\end{proof}

\subsection{The Maurer--Cartan form and an intrinsic differential of the exponential}
None of the assertions on the page requires a matrix embedding. For an
arbitrary finite-dimensional real or complex Lie group $G$ with Lie algebra
$\g$, define the left Maurer--Cartan form by
$\theta_g=(L_{g^{-1}})_*:T_gG\to T_eG=\g$; in a matrix group it is
$g^{-1}\dd g$. On a left-invariant field it takes the constant value
$\theta(U^L)=U$.

\begin{proposition}[Maurer--Cartan equation]\label{prop:MC}
The $\g$-valued one-form $\theta$ satisfies
\begin{equation}\label{eq:MC}
 \dd\theta+\tfrac12[\theta,\theta]=0,\qquad
 \tfrac12[\theta,\theta](U,V):=[\theta(U),\theta(V)].
\end{equation}
Consequently, for a smooth two-parameter map $g(s,t)$ into $G$, the
left-trivialized partial derivatives
$\alpha=\theta(\partial_sg)$ and $\eta=\theta(\partial_tg)$ satisfy the
mixed-variation identity
\begin{equation}\label{eq:mixedvariation}
 \partial_s\eta-\partial_t\alpha+[\alpha,\eta]=0.
\end{equation}
\end{proposition}
\begin{proof}
For a one-form $\omega$ and vector fields $U,V$ one has
$\dd\omega(U,V)=U(\omega(V))-V(\omega(U))-\omega([U,V])$. Evaluate on
left-invariant fields $U^L,V^L$: the functions $\theta(V^L)=V$ and
$\theta(U^L)=U$ are constant, so the first two terms vanish, and
$[U^L,V^L]=[U,V]^L$ by definition of the Lie algebra bracket; hence
$\dd\theta(U^L,V^L)=-\theta([U,V]^L)=-[U,V]$. On the other hand
$\frac12[\theta,\theta](U^L,V^L)=[\theta(U^L),\theta(V^L)]=[U,V]$. Both
sides of \eqref{eq:MC} are tensorial (bilinear over functions) and
left-invariant fields span every tangent space, so \eqref{eq:MC} holds
identically. In a matrix group the same result is the direct computation
$\dd(g^{-1})=-g^{-1}(\dd g)g^{-1}$, whence
$\dd(g^{-1}\dd g)=-g^{-1}\dd g\wedge g^{-1}\dd g$.

For \eqref{eq:mixedvariation}, pull $\theta$ back by $g$:
$g^*\theta=\alpha\,\dd s+\eta\,\dd t$. Pullback commutes with $\dd$ and
with the bracket of forms, so \eqref{eq:MC} gives
$\dd(g^*\theta)=-[\alpha,\eta]\,\dd s\wedge\dd t$, while directly
$\dd(\alpha\,\dd s+\eta\,\dd t)=(\partial_s\eta-\partial_t\alpha)\,\dd s\wedge\dd t$.
Comparing coefficients gives \eqref{eq:mixedvariation}. (In a matrix group:
$\partial_s(g^{-1}\partial_tg)=-\alpha\eta+g^{-1}\partial_s\partial_tg$ and
$\partial_t(g^{-1}\partial_sg)=-\eta\alpha+g^{-1}\partial_t\partial_sg$;
subtracting gives $-[\alpha,\eta]$.)
\end{proof}

\begin{proposition}[Intrinsic differential of the exponential]\label{prop:intrinsicdexp}
For $X,H\in\g$,
\begin{equation}\label{eq:intrinsicdexp}
 (L_{\exp(-X)})_*(\dd\exp)_X(H)=\phi(\ad_X)H=\int_0^1e^{-u\ad_X}H\,\dd u,
\end{equation}
and similarly $(R_{\exp(-X)})_*(\dd\exp)_X(H)=\phi(-\ad_X)H$. In
particular $(\dd\exp)_0=\id_\g$, and $\Ad_{\exp X}=e^{\ad_X}$ intrinsically.
\end{proposition}
\begin{proof}
Put $g(s,t)=\exp\bigl(s(X+tH)\bigr)$. For fixed $t$, $s\mapsto g(s,t)$ is
the one-parameter subgroup with velocity $X+tH$, so its left-trivialized
velocity is $\alpha(s,t)=\theta(\partial_sg)=X+tH$. Let
$\eta(s,t)=\theta(\partial_tg)$; then $\eta(0,t)=0$ because $g(0,t)=e$
for all $t$. By \eqref{eq:mixedvariation},
$\partial_s\eta=\partial_t\alpha-[\alpha,\eta]=H-[X+tH,\eta]$. At $t=0$
this is the linear equation $\partial_s\eta=H-\ad_X\eta$ with
$\eta(0)=0$, whose unique solution is
$\eta(s)=\int_0^se^{-(s-u)\ad_X}H\,\dd u$ (check by differentiation:
$\partial_s\eta=H-\ad_X\eta$). At $s=1$, substituting $v=1-u$,
$\eta(1,0)=\int_0^1e^{-v\ad_X}H\,\dd v=\phi(\ad_X)H$ by the termwise
integration of \cref{thm:dexp}. Since $\eta(1,0)=\theta(\partial_tg(1,0))
=(L_{\exp(-X)})_*(\dd\exp)_X(H)$, this is \eqref{eq:intrinsicdexp}. The
right-trivialized version follows by applying the same argument to the
right Maurer--Cartan form, or from
$(R_{\exp(-X)})_*=\Ad_{\exp X}\circ(L_{\exp(-X)})_*$ together with
$\Ad_{\exp X}\phi(\ad_X)=e^{\ad_X}\phi(\ad_X)=\phi(-\ad_X)$. At $X=0$ the
formula gives $(\dd\exp)_0(H)=H$. Finally, differentiating the
one-parameter group of automorphisms $t\mapsto\Ad_{\exp(tX)}$ gives
$\frac{\dd}{\dd t}\Ad_{\exp(tX)}=\ad_X\Ad_{\exp(tX)}$ with
$\Ad_{\exp(0)}=I$, whose unique solution is $e^{t\ad_X}$; this is the
intrinsic Campbell identity.
\end{proof}

\subsection{Local exponential coordinates and the local BCH law}
\begin{theorem}[Local BCH multiplication]\label{thm:local}
Let $G$ be a finite-dimensional real or complex Lie group. There is a
neighborhood $U$ of $0$ in $\g$ on which $\exp$ is a diffeomorphism onto a
neighborhood of the identity. After shrinking $U$, for all $X,Y\in U$,
\begin{equation}\label{eq:localproduct}
 \exp X\exp Y=\exp\bigl(\BCH(X,Y)\bigr),
\end{equation}
where the universal series converges absolutely in $\g$. Hence the local
multiplication law of $G$ is determined entirely by the Lie bracket.
\end{theorem}
\begin{proof}
By \cref{prop:intrinsicdexp}, $(\dd\exp)_0=\id$, so the inverse function
theorem provides the chart. Choose a norm on $\g$ and $\kappa$ with
$\norm{[u,v]}\leq\kappa\norm u\norm v$. By \cref{thm:convergence}(iv), for
$\kappa(\norm X+\norm Y)<\log2$ the series $Z(t)=\BCH(X,tY)=\sum_nZ_n(X,tY)$
converges absolutely, uniformly for $t\in[0,1]$, and so does the series
with $X,Y$ replaced by $rX,rY$ for some $r>1$, provided $X,Y$ are small
enough. Consequently $Z(t)$ is given by a power series in $t$ with radius
greater than $1$, and it may be differentiated termwise in $t$. The
differential equation \eqref{eq:ode}, $Z'(t)=\beta(\ad_{Z(t)})Y$, holds in
the formal setting coefficient by coefficient in $t$ (each coefficient
being a Lie polynomial identity in $X,Y$, valid in $\g$ by universality,
\cref{sec:universal}); for $\norm{Z(t)}$ small enough that
$\norm{\ad_{Z(t)}}<2\pi$, the right side is a convergent series in $t$ with
the same coefficients, so the equation holds analytically.

Now let $L(t)=\log(\exp X\exp(tY))$ in the exponential chart, defined for
small $X,Y$ and $t\in[0,1]$ after shrinking $U$ so that the relevant
products stay in the chart. Its left-trivialized derivative is
$\theta\bigl(\tfrac{\dd}{\dd t}\exp X\exp(tY)\bigr)=Y$ (the curve
$t\mapsto\exp X\exp(tY)$ is the left translate by $\exp X$ of a
one-parameter subgroup). By \cref{prop:intrinsicdexp} the same
left-trivialized derivative equals $\phi(\ad_{L(t)})L'(t)$. For $L$ small,
$\phi(\ad_L)$ is invertible with inverse $\beta(\ad_L)$, so
$L'=\beta(\ad_L)Y$ with $L(0)=X$. The map $L\mapsto\beta(\ad_L)Y$ is
analytic on a neighborhood of $0$, so the initial value problem has a
unique solution by the Picard--Lindel\"of theorem. Both $L(t)$ and $Z(t)$
solve it; hence $L(1)=Z(1)=\BCH(X,Y)$, which is \eqref{eq:localproduct}
since $\exp$ is injective on $U$.
\end{proof}
The local assertion is not a global classification. The additive group
$\R$ and the circle group $S^1=\{e^{\ii t}\}$ have the same one-dimensional
abelian Lie algebra, but one is noncompact and the other compact, so they
are not isomorphic as Lie groups; they differ in topology, not in local
structure. The identity map of their Lie algebras integrates from $\R$ to
$S^1$ as $t\mapsto e^{\ii t}$, but not from $S^1$ to $\R$: a continuous
homomorphism from the compact group $S^1$ to $\R$ has compact image, which is
a compact subgroup of $\R$, hence $\{0\}$, so its differential is zero.

\subsection{Homomorphisms, representations, and integration}
\begin{theorem}[Local Lie correspondence]\label{thm:localcorr}
Let $f:\g\to\h$ be a homomorphism of Lie algebras of Lie groups $G,H$. Then
\begin{equation}\label{eq:localhom}
 F(\exp_GX)=\exp_H(fX)
\end{equation}
defines a smooth map on an identity neighborhood of $G$ which is a local
homomorphism: $F(gh)=F(g)F(h)$ whenever $g,h,gh$ lie in a suitable smaller
neighborhood. Conversely, the differential at the identity of a smooth
local homomorphism preserves Lie brackets, and a local homomorphism is
determined near the identity by that differential.
\end{theorem}
\begin{proof}
On the exponential neighborhood of \cref{thm:local}, \eqref{eq:localhom}
defines a smooth map, being the composition $\exp_H\circ f\circ\log_G$. By
\cref{prop:symmetry}, $f$ preserves every $Z_n$, and by continuity it
preserves the convergent sum: $f(\BCH(X,Y))=\BCH(fX,fY)$. Hence, for $X,Y$
small enough that \eqref{eq:localproduct} applies in both groups,
\begin{align*}
 F(\exp X\exp Y)&=F\bigl(\exp\BCH(X,Y)\bigr)=\exp_H\bigl(f\,\BCH(X,Y)\bigr)
 =\exp_H\bigl(\BCH(fX,fY)\bigr)\\
 &=\exp_H(fX)\exp_H(fY)=F(\exp X)F(\exp Y).
\end{align*}
For the converse, let $F$ be a smooth local homomorphism with differential
$A=(\dd F)_e$. Differentiating $F(g\exp(tX))=F(g)F(\exp(tX))$ at $t=0$
shows that $F$ relates the left-invariant field $X^L$ on $G$ to the field
$(AX)^L$ on $H$, wherever both are defined. Related vector fields have
related brackets (a standard consequence of the chain rule applied to
$F_*$), so $[X^L,Y^L]=[X,Y]^L$ is related to $[(AX)^L,(AY)^L]=[AX,AY]^L$;
evaluating at $e$ gives $A[X,Y]=[AX,AY]$. Finally, since $F$ maps the
one-parameter subgroup $\exp(tX)$ to a local one-parameter subgroup of $H$
with initial velocity $AX$, uniqueness of integral curves gives
$F(\exp(tX))=\exp_H(tAX)$ for small $t$; as $\exp$ is onto an identity
neighborhood, $F$ is determined near $e$ by $A$.
\end{proof}

\begin{theorem}[Integration from a simply connected group]\label{thm:integration}
If $G$ is connected and simply connected and $H$ is a Lie group, every Lie
algebra homomorphism $f:\g\to\h$ integrates to a unique smooth homomorphism
$F:G\to H$ with $(\dd F)_e=f$. In particular every finite-dimensional
representation of $\g$ integrates to a representation of $G$.
\end{theorem}
\begin{proof}
Let $F_0$ be the local homomorphism of \cref{thm:localcorr}, defined and
multiplicative on an open identity neighborhood $W$ in the sense that
$F_0(uv)=F_0(u)F_0(v)$ whenever $u,v,uv\in W$. Choose a smaller symmetric
identity neighborhood $W_1$ with $W_1W_1\subseteq W$.

\emph{Step 1: paths.} Let $\gamma:[0,1]\to G$ be a continuous path with
$\gamma(0)=e$ and $\gamma(1)=g$. By uniform continuity there is a partition
$0=t_0<t_1<\cdots<t_m=1$ such that every increment
$u_j=\gamma(t_{j-1})^{-1}\gamma(t_j)$ lies in $W_1$. Define
\[
 F_\gamma(g)=F_0(u_1)F_0(u_2)\cdots F_0(u_m).
\]
This does not depend on the partition. Inserting one additional point $t'$
between $t_{j-1}$ and $t_j$ replaces $u_j$ by the product $u'u''$ of two
increments; if the partition is fine enough that $u',u''\in W_1$, then
$u'u''=u_j\in W$ and $F_0(u_j)=F_0(u')F_0(u'')$ by local multiplicativity,
so the product is unchanged. Any two sufficiently fine partitions have a
common refinement obtained by inserting points one at a time, so
$F_\gamma(g)$ is well defined.

\emph{Step 2: homotopy invariance.} Let $\Gamma:[0,1]^2\to G$ be a
continuous homotopy with $\Gamma(\cdot,0)=\gamma_0$, $\Gamma(\cdot,1)=\gamma_1$,
$\Gamma(0,\cdot)=e$, $\Gamma(1,\cdot)=g$. By uniform continuity there is a
rectangular grid so fine that for any two grid points $p,q$ of the same
closed cell, $\Gamma(p)^{-1}\Gamma(q)\in W_1$. Around a single cell with
corners $a,b,c,d$ (in cyclic order), the four increments
$u=\Gamma(a)^{-1}\Gamma(b)$, $v=\Gamma(b)^{-1}\Gamma(c)$,
$u'=\Gamma(a)^{-1}\Gamma(d)$, $v'=\Gamma(d)^{-1}\Gamma(c)$ satisfy
$uv=u'v'=\Gamma(a)^{-1}\Gamma(c)\in W_1$, so local multiplicativity gives
$F_0(u)F_0(v)=F_0(uv)=F_0(u')F_0(v')$. Hence the products of $F_0$ along the
two edge paths of the cell from $a$ to $c$ coincide. Sweeping across the
grid one cell at a time, replacing the lower-right edge pair of a cell by
its upper-left pair, converts the grid path along the bottom and right
boundary into the path along the left and top boundary without changing the
product. The left and right boundary edges are constant paths (at $e$ and
at $g$) and contribute factors $F_0(e)=e$. The bottom and top boundary
paths are $\gamma_0$ and $\gamma_1$ evaluated at the grid points, whose
increments form fine partitions in the sense of Step 1. Therefore
$F_{\gamma_0}(g)=F_{\gamma_1}(g)$.

\emph{Step 3: definition of $F$.} Since $G$ is simply connected, any two
paths from $e$ to $g$ are homotopic with fixed endpoints, and since $G$ is
connected such paths exist. Define $F(g)=F_\gamma(g)$ for any such path.

\emph{Step 4: multiplicativity, smoothness, differential.} Given paths
$\gamma$ from $e$ to $g$ and $\delta$ from $e$ to $h$, the concatenation of
$\gamma$ with the left translate $g\delta$ is a path from $e$ to $gh$. The
increments of $g\delta$ are $(g\delta(t_{j-1}))^{-1}g\delta(t_j)
=\delta(t_{j-1})^{-1}\delta(t_j)$, the original increments of $\delta$.
Hence $F(gh)=F(g)F(h)$. For $u\in W_1$, appending to a path to $g$ the
short path $t\mapsto g\exp(t\log u)$ from $g$ to $gu$ gives
$F(gu)=F(g)F_0(u)$; thus $F$ agrees near $g$ with the smooth map
$u\mapsto F(g)F_0(u)$, so $F$ is smooth, and at $g=e$ it agrees with $F_0$,
so $(\dd F)_e=f$.

\emph{Step 5: uniqueness.} Any homomorphism $F'$ with differential $f$
satisfies $F'(\exp X)=\exp_H(fX)$ (it maps one-parameter subgroups to
one-parameter subgroups with the given initial velocity), so $F'=F_0=F$ on
the exponential neighborhood. The set where $F'=F$ is a subgroup
containing an identity neighborhood, hence an open subgroup; its cosets
are open, so it is also closed; connectedness makes it all of $G$.

Taking $H=\GL(E)$ for a finite-dimensional vector space $E$ gives the
statement about representations.
\end{proof}
Without simple connectedness the path construction can yield a
nontrivial value around a loop; this is the precise global obstruction
absent from the local BCH law. Descent from a simply connected covering
group to a specified quotient requires the covering kernel to act trivially.

\subsection{Nilpotent Lie algebras and the polynomial group law}\label{sec:nilpotent}
Write $\g^1=\g$ and $\g^{r+1}=[\g,\g^r]$ for the lower central series. The
algebra is nilpotent of class at most $c$ if $\g^{c+1}=0$.
\begin{theorem}[Termination and the nilpotent group law]\label{thm:nilpotent}
In a nilpotent Lie algebra of class at most $c$, every $Z_n$ with $n>c$
vanishes identically, and
\begin{equation}\label{eq:nilpotentlaw}
 X*Y=\BCH(X,Y)=\sum_{n=1}^{c}Z_n(X,Y)
\end{equation}
is a polynomial group law on the vector space $\g$ with identity $0$ and
inverse $-X$. For finite-dimensional real $\g$ this makes $\g$ a connected,
simply connected Lie group whose Lie algebra is $\g$ with its original
bracket and whose exponential map is the identity of $\g$. Any connected
simply connected Lie group with Lie algebra $\g$ is isomorphic to it, and
in every Lie group with Lie algebra $\g$ the identity
$\exp X\exp Y=\exp\BCH(X,Y)$ holds for all $X,Y\in\g$.
\end{theorem}
\begin{proof}
By the Jacobi identity in the form \eqref{eq:derivation}, induction on $r$
gives $[\g^r,\g^s]\subseteq\g^{r+s}$: for $r=1$ this is the definition,
and if it holds for $r$ then for $x\in\g$, $y\in\g^r$, $z\in\g^s$,
$[[x,y],z]=[x,[y,z]]-[y,[x,z]]\in[\g,\g^{r+s}]+[\g^r,\g^{s+1}]\subseteq\g^{r+s+1}$.
Hence every bracket monomial of degree $n$ in elements of $\g$ lies in
$\g^n$, and vanishes for $n>c$. Every $Z_n$ is such a combination, so
$Z_n=0$ for $n>c$.

The formal identities \eqref{eq:assoc} and \eqref{eq:unitinverse} are
identities between Lie series in three free variables. Evaluating them in
$\g$ turns both sides into finite polynomial identities, because all Lie
monomials of degree greater than $c$ vanish, including those produced by
substituting one polynomial into another. Thus
$(X*Y)*T=X*(Y*T)$, $X*0=0*X=X$, and $X*(-X)=0$ hold for all $X,Y,T\in\g$,
not merely near zero. So $(\g,*)$ is a group, with polynomial multiplication
and inversion; it is a Lie group whose underlying manifold is a vector
space, hence connected and simply connected.

Its Lie algebra is computed from the group commutator. Write $a=sX$,
$b=tY$, and let $O(3)$ denote terms of total degree at least three in
$s,t$. Using $u*v=u+v+\frac12[u,v]+O(3)$ three times,
\begin{align*}
 (a*b)*(-a)&=\bigl(a+b+\tfrac12[a,b]\bigr)-a+\tfrac12[a+b,-a]+O(3)
 =b+[a,b]+O(3),\\
 \bigl((a*b)*(-a)\bigr)*(-b)&=\bigl(b+[a,b]\bigr)-b+\tfrac12[b,-b]+O(3)
 =[a,b]+O(3)=st[X,Y]+O(3).
\end{align*}
The mixed second derivative $\partial_s\partial_t$ at $s=t=0$ of the group
commutator $g_sh_tg_s^{-1}h_t^{-1}$ of two one-parameter subgroups is the
Lie bracket of their initial velocities in the Lie algebra of any Lie
group, so the bracket of $(\g,*)$ is the original bracket of $\g$.
(Alternatively, the left-invariant field of $X$ is $g\mapsto\frac{\dd}{\dd t}
g*(tX)|_{t=0}$, and computing brackets of these polynomial fields gives the
same result.) Since $Z(sX,tX)=(s+t)X$ (all brackets of multiples of $X$
vanish), the one-parameter subgroups are $t\mapsto tX$, so the exponential
map is the identity.

If $G$ is a connected simply connected Lie group with Lie algebra $\g$,
\cref{thm:integration} integrates the identity map of $\g$ in both
directions, giving homomorphisms $(\g,*)\to G$ and $G\to(\g,*)$ whose
composites have differential the identity, hence are the identity by the
uniqueness part; so they are inverse isomorphisms. In particular the
exponential map of $G$ is a global diffeomorphism. For an arbitrary Lie
group $G'$ with Lie algebra $\g$, the homomorphism $(\g,*)\to G'$
integrating the identity map sends $X$ to $\exp_{G'}X$ (it maps
one-parameter subgroups to one-parameter subgroups), and multiplicativity
reads $\exp_{G'}(X*Y)=\exp_{G'}X\exp_{G'}Y$.
\end{proof}
Termination is a statement about Lie brackets, not about matrix powers: a
nilpotent Lie algebra may be represented by matrices that are not
themselves nilpotent.

For a basis $P,Q,C$ with $[P,Q]=C$ and all other brackets zero (the
Heisenberg algebra, of class two), \eqref{eq:nilpotentlaw} reads
\begin{equation}\label{eq:heisenberglaw}
 (x,y,z)*(x',y',z')=\bigl(x+x',\,y+y',\,z+z'+\tfrac12(xy'-yx')\bigr)
\end{equation}
in coordinates $xP+yQ+zC$, since $\frac12[xP+yQ,x'P+y'Q]=\frac12(xy'-yx')C$.
A concrete matrix instance is $P=E_{12}$, $Q=E_{23}$, $C=E_{13}$ in the
strictly upper triangular $3\times3$ matrices, where both sides of
$e^Xe^Y=e^{X+Y+[X,Y]/2}$ are finite matrix polynomials. This group law
underlies the Weyl and displacement identities of \cref{sec:quantum}.
